\documentclass[11pt]{article}

\usepackage[a4paper,margin=29mm]{geometry}
\usepackage{amsmath,amssymb,amsthm,mathtools}
\usepackage{array,booktabs}
\usepackage{microtype}
\usepackage[numbers,sort&compress]{natbib}
\usepackage[hidelinks]{hyperref}

\newtheorem{theorem}{Theorem}
\newtheorem{proposition}[theorem]{Proposition}
\newtheorem{corollary}[theorem]{Corollary}
\newtheorem{lemma}[theorem]{Lemma}
\theoremstyle{definition}
\newtheorem{conjecture}[theorem]{Conjecture}
\newtheorem{remark}[theorem]{Remark}
\newtheorem{observation}[theorem]{Observation}

\newcommand{\R}{\mathbb{R}}
\newcommand{\disc}{\Delta}
\newcommand{\eps}{\varepsilon}
\newcommand{\W}{W}

\title{Detector Chains and a Lambert-$W$ Lower Bound\\
for Machine-Dependent Chairman Assignment}
\author{Angus Muffatti}
\date{25 July 2026 --- draft}

\begin{document}
\maketitle

\begin{abstract}
For fractional assignments with machine-dependent weights $d_{ij}>0$, let
$\kappa$ denote the supremum, over all finite instances, of the minimum over
integral assignments of the maximum row-and-prefix discrepancy divided by
$D=\max_{i,j}d_{ij}$.  A companion paper refuted the conjecture $\kappa\le1$
of Liu and Reis and proved $\kappa\ge4-2\sqrt2$ using a three-column
detector gadget.  Here we replace the three-column detector by a
\emph{chain} of $k-1$ detector columns and optimize the resulting family
exactly.  The balanced chain design has the closed-form value
\[
 \tau(k)\;=\;\max_{d}\;
 \frac{d\bigl(k-1-(k-2)\,d^{1/(k-2)}\bigr)+1}{1+d}
\]
(the maximum over weights keeping the design feasible), and
\[
 \lim_{k\to\infty}\tau(k)\;=\;1+\W(1/e)\;\approx\;1.278465,
\]
where $\W$ is the Lambert $\W$ function; at the optimum the pleasant identity
$\tau=1+d^*$ holds, with $\ln d^*=-(1+d^*)$.  Consequently
$\kappa\ge1+\W(1/e)$.  Every rung is certified in exact rational
arithmetic: we exhibit explicit instances with certified prefix
discrepancy $1.2727$ (depth $40$), and an independent mixed-integer check
confirms a concrete $119\times236$ instance beating $7/6$.  We further
report an extensive, exactly verified experimental study of mechanisms
beyond chains --- banked helpers, bank-manufacturing ``shaper'' blocks, and
block games with forced values above $3/2$ at stated bank states --- all of
which fail, measurably and instructively, to compose into instances above
the chain limit.  We conjecture that $\kappa$ is finite and discuss a
scheduler-side program for a matching upper bound, as well as the bounded
weight-ratio parameterization $\kappa(\rho)$.
\end{abstract}

\section{Introduction}

Let $x\in[0,1]^{m\times n}$ be a fractional assignment,
$\sum_ix_{ij}=1$ for every column $j$, and let $d\in\R_{>0}^{m\times n}$ be
a matrix of machine-dependent weights.  For an integral assignment
$y\in\{0,1\}^{m\times n}$ with the same column sums, write
\[
 \disc_i(t)=\sum_{j=1}^{t}d_{ij}(x_{ij}-y_{ij})
\]
for the row-prefix discrepancy.  Liu and Reis \citep{LiuReis2026}, having
proved the common-weight prefix guarantee that resolves a special case of
the Morell--Skutella conjecture \citep{MorellSkutella2022}, conjectured
that machine-dependent weights always admit $y$ with
$|\disc_i(t)|\le D:=\max_{i,j}d_{ij}$ (their Conjecture 19; Conjecture 3 in
the arXiv version).  The companion paper \citep{Muffatti2026} refuted this
with an $11\times15$ rational instance and strengthened the refutation to a
quantitative lower bound: writing
\[
 \kappa\;=\;\sup_{\text{instances}}\;
 \min_{y}\;\max_{i,t}\;\frac{|\disc_i(t)|}{D},
\]
it proved $\kappa\ge4-2\sqrt2\approx1.1716$, via a three-column
\emph{detector gadget} whose split was optimized, and promised a
systematic treatment of deeper detectors.  This paper is that treatment.

\paragraph{Results.}
Our main construction generalizes the three-column detector to a
\emph{detector chain} of depth $k$: one hypothesis column followed by
$k-1$ chained detector columns on a fresh row pair, all coupled to a
global accumulator row.  Optimizing the weights and splits exactly yields
three statements.

\begin{enumerate}
\item \textbf{A closed form per depth} (Theorem~\ref{thm:tau}).  The
balanced chain of depth $k$ forces prefix discrepancy arbitrarily close to
\[
 \tau(k)=\max_{d\ \text{admissible}}
 \frac{d\bigl(k-1-(k-2)d^{1/(k-2)}\bigr)+1}{1+d},
\]
where admissibility means the balance splits lie in $[0,1]$
(equivalently the displayed value is at most $1+d$), and
$\tau(3)=4-2\sqrt2$ recovers the companion paper's constant.
\item \textbf{A Lambert-$W$ limit} (Theorem~\ref{thm:limit}).
$\tau(k)$ converges to $1+\W(1/e)$, where $d^*=\W(1/e)\approx0.27846$
satisfies $\ln d^*=-(1+d^*)$ and the limiting optimal value is exactly
$1+d^*$.  Hence $\kappa\ge1+\W(1/e)$.
\item \textbf{Exact certificates} (Section~\ref{sec:certificates}).
Every claim is machine-verified in exact rational arithmetic.  Concrete
certified instances include a depth-$4$ family member
($\tau=103/85$) whose $119\times236$ instance provably exceeds $7/6$
(independent mixed-integer infeasibility check) and a depth-$40$ instance
with certified discrepancy $\ge1.272743$, within $0.006$ of the limit.
\end{enumerate}

The second half of the paper (Section~\ref{sec:beyond}) reports what lies
beyond chains.  It is experimental in the strict sense that every stated
number is an exactly verified fact about a finite object, while the
section's larger claims are explicitly conjectural.  In brief: block-level
mechanisms exist whose forced values exceed $3/2$ \emph{given} favorable
bank states, and ``shaper'' blocks exist that manufacture such states
against every legal schedule at fixed budgets; yet every attempt to
compose these components into an actual instance has failed, for reasons
we can now name and measure (entry pollution, accumulator venting, and the
failure of conjunctive bank pinning).  The evidence points to a
\emph{composition barrier}: the constant $\kappa$ appears to be governed
not by the strength of local gadgets, which is unbounded, but by a
corridor-management game that the scheduler keeps winning.  We conjecture
$\kappa<\infty$ and outline the upper-bound program this suggests.

\paragraph{Why care.}
Beyond settling the strength of the natural gadget family, the
parameterized reading is relevant to scheduling practice: $\kappa$ prices
\emph{unstructured heterogeneity} in prefix-fair scheduling.
Section~\ref{sec:rho} records the bounded-ratio corollary: with all
weights in $[D/\rho, D]$, the constructions degrade gracefully,
$\kappa(\rho)\to1+\W(1/e)$ as $\rho\to\infty$, while $\kappa(1)=1$ by the
common-weight theorem; the curve between is open and is, in our view, the
right target for applications.

\section{The chain gadget}\label{sec:gadget}

Fix a depth $k\ge3$, a detector share $p\in(0,1)$, an off-support
weight $\eta>0$, and a block count $N$.  The rows are
$B,A_1,H_1,\ldots,A_N,H_N$; the columns come in $N$ consecutive blocks
\[
 J_b,\;C^b_2,\;C^b_3,\;\ldots,\;C^b_k\qquad(b=1,\ldots,N).
\]
Within block $b$ (write $A=A_b$, $H=H_b$), the nonzero fractional entries
and exceptional weights are:
\[
\begin{array}{c|cc|cc}
 \text{column} & x_A & x_{\cdot} & d_A & d_{\cdot}\\
\hline
 J   & 1-p & p\ (B) & 1 & 1\\
 C_j & x_j & 1-x_j\ (H) & d_j & 1\qquad(j=2,\ldots,k)
\end{array}
\]
with weights $d_2\le\cdots\le d_{k-1}\le d_k=1$ and splits
$x_j\in[0,1]$ to be chosen; every unlisted weight equals $\eta$.  All
weights are strictly positive and $D=1$.

The forcing logic generalizes the three-column case.  Suppose a
hypothetical assignment keeps all prefixes within a budget $c$, and
suppose it denies $J$ to $A$.  Then $A$ sits near $1-p$, and each chain
column presents the scheduler with a dilemma: assigning $C_j$ to $A$
relieves $A$ by $d_j(1-x_j)$ but denies $H$ its share $1-x_j$, while any
other choice denies $A$ the amount $d_jx_j$ and fires immediately.  The
weights grow along the chain, so the relief purchased early is small
while the punishment for late deviation stays high; the terminal column
finally forces a violation on one of the two rows.

\begin{lemma}[block heights]\label{lem:heights}
Deny $J$ to $A$ and let the scheduler play the block columns arbitrarily.
With the convention that sums over $i$ range over chain columns before the
current one, every continuation attains one of the following prefix
heights on $A$ or $H$:
\begin{align*}
 \mathrm{dev}_j&=(1-p)-\textstyle\sum_{i<j}d_i(1-x_i)+d_jx_j
 &&(\text{first deviation at }C_j,\ j<k),\\
 \mathrm{term}_A&=(1-p)-\textstyle\sum_{i<k}d_i(1-x_i)+x_k
 &&(\text{terminal denied to }A),\\
 \mathrm{term}_H&=\textstyle\sum_{i<k}(1-x_i)+(1-x_k)
 &&(\text{full compliance}).
\end{align*}
These heights are exact up to the $\eta$-bookkeeping of
Lemma~\ref{lem:budget}: a column sent to neither support row denies both,
so it only raises the listed heights.
\end{lemma}

\begin{proof}
If the first chain column not assigned to $A$ is $C_j$ with $j<k$, then at
that prefix $A$ has received exactly the earlier chain columns and been
denied $J$'s share and $C_j$'s share, giving $\mathrm{dev}_j$; this is
attained at the prefix of $C_j$ itself, before any later
column can repair it.  If all of $C_2,\ldots,C_{k-1}$ go to $A$ and $C_k$
does not, $A$'s height at the terminal prefix is $\mathrm{term}_A$.  If
every chain column goes to $A$, then $H$ has been denied its share
$1-x_j$ of every chain column, and its height at the terminal prefix is
$\mathrm{term}_H$.  Assignments off the support deny both rows and can
only increase these values (they subtract nothing and cost the recipient
only $\eta$, which Lemma~\ref{lem:budget} accounts).
\end{proof}

\begin{lemma}[generalized budget lemma]\label{lem:budget}
Let $\tau$ be a value such that the $p=0$ idealizations of the heights of
Lemma~\ref{lem:heights} (the same expressions with leading term $1$ in
place of $1-p$) are all at least $\tau$.  Then every integral assignment
of the $N$-block instance has some row-prefix discrepancy at least
\[
 \mathrm{LB}\;=\;\min\Bigl\{\ \tau-p-\bigl(k(N-1)+1\bigr)\eta,\quad
 Np-N(k-1)\eta\ \Bigr\}.
\]
\end{lemma}

\begin{proof}
The fresh rows of block $b$ appear off support in at most $k(N-1)$ earlier
columns, each costing $\eta$, so they enter the block no lower than
$-k(N-1)\eta$.  Within the block, $A$ is on support in every column, so it
takes no further $\eta$ hits; $H$ can additionally absorb $J$ off support
once.  The heights of Lemma~\ref{lem:heights} lose exactly $p$ on the
$A$-side entries (their first term is $1-p$ rather than $1$) and nothing
on $\mathrm{term}_H$; combining, every height is at least
$\tau-p-(k(N-1)+1)\eta$.  If this exceeds the running budget, denying any
$J_b$ is impossible, so all $N$ detectors are assigned to their $A_b$;
then $B$ is denied its share $p$ in every block and can be relieved only
by off-support assignments of chain columns, each worth $\eta$, of which
there are at most $N(k-1)$.  At the final prefix
$\disc_B\ge Np-N(k-1)\eta$.  Applying the argument with the budget equal
to the assignment's actual maximum prefix discrepancy yields the claim.
\end{proof}

\section{Optimizing the chain}\label{sec:optimize}

Equalize the heights.  In the idealization $p\to0$, setting all
$\mathrm{dev}_j$, $\mathrm{term}_A$, $\mathrm{term}_H$ equal to a common
value $\tau$ forces the \emph{balance solution}
\begin{equation}\label{eq:balance}
 x_{j+1}=\frac{d_j}{d_{j+1}}\quad(2\le j\le k-1,\ d_k:=1),
 \qquad
 x_2=\frac{\tau-1}{d_2},
\end{equation}
as one checks by subtracting consecutive height expressions.  Substituting
into $\mathrm{term}_H=\tau$ and writing
\[
 R\;=\;\sum_{j=2}^{k-1}\frac{d_j}{d_{j+1}}
\]
gives $(k-1)-\frac{\tau-1}{d_2}-R=\tau$, that is,
\begin{equation}\label{eq:tau}
 \tau\;=\;\frac{d_2\,(k-1-R)+1}{1+d_2}.
\end{equation}

\begin{lemma}[geometric weights are optimal]\label{lem:amgm}
For fixed $d_2$, the value \eqref{eq:tau} is maximized when
$d_2,\ldots,d_{k-1},1$ is a geometric progression, in which case
$R=(k-2)\,d_2^{1/(k-2)}$.
\end{lemma}

\begin{proof}
The ratios $r_j=d_j/d_{j+1}$ satisfy $\prod_{j=2}^{k-1}r_j=d_2/d_k=d_2$.
By AM--GM, $R=\sum r_j\ge(k-2)\bigl(\prod r_j\bigr)^{1/(k-2)}
=(k-2)d_2^{1/(k-2)}$, with equality iff all ratios are equal.  Since
\eqref{eq:tau} is decreasing in $R$, the geometric choice is optimal.
\end{proof}

Two splits of the balance solution can leave the unit box, so record the
condition under which the design is realizable.  Call a weight
$d\in(0,1]$ \emph{admissible} for depth $k$ if the geometric balance
design has all splits in $[0,1]$.  The interior ratios equal
$d^{1/(k-2)}\le1$ automatically, so admissibility is exactly the condition
$x_2=(\tau-1)/d\le1$, that is, $\tau(k,d)\le1+d$, where
\[
 \tau(k,d)\;:=\;\frac{d\bigl(k-1-(k-2)\,d^{1/(k-2)}\bigr)+1}{1+d}.
\]
(Small $d$ can be inadmissible at large $k$; every table point of
Section~\ref{sec:certificates} is verified admissible.)

Equalizing is also not an arbitrary choice: the terminal helper height
binds at every optimum except in one boundary configuration, which we
isolate because a naive monotonicity argument overlooks it.

\begin{lemma}\label{lem:tight}
Fix the weights $d$ and, for splits $x\in[0,1]^{k-1}$, put
\[
 f(x)=\min\{\mathrm{dev}_2,\ldots,\mathrm{dev}_{k-1},\mathrm{term}_A\},
 \qquad
 g(x)=\mathrm{term}_H
\]
in the $p=0$ idealization.  At every maximizer of
$\min\{f,g\}$, either $f(x)=g(x)$, or else $f(x)<g(x)$ and $x_2=1$, in
which case $\min\{f,g\}=1+d_2$.
\end{lemma}

\begin{proof}
Each $A$-side height is nondecreasing in every coordinate: $\mathrm{dev}_j$
involves $x_i$ ($i<j$) only through $+d_ix_i$ and $x_j$ through $+d_jx_j$,
and $\mathrm{term}_A$ likewise has nonnegative coefficients.  In contrast
$g=\sum_{i=2}^{k-1}(1-x_i)+(1-x_k)$ is strictly decreasing in every
coordinate.

Suppose $f(x)>g(x)$, so $\min\{f,g\}=g$.  At $x=0$ we have $g=k-1\ge2$
while $f\le\mathrm{dev}_2=1$, so this case forces $x\ne0$.  Decreasing a
positive coordinate slightly raises $g$ strictly and cannot raise $f$, and
$f>g$ persists for a small enough step, so $\min\{f,g\}$ increases,
contradicting optimality.

Now suppose $g(x)>f(x)$, so $\min\{f,g\}=f$, and let $B$ be the set of
$A$-side heights attaining $f$.  Raising every split that is below $1$ by
a common small $\delta$ strictly increases any height in $B$ that has at
least one split below $1$, while $g$ decreases but retains its slack; so
if every member of $B$ had a split below $1$, then $f$ would strictly
increase, a contradiction.  Hence some member of $B$ has all of its splits
equal to $1$.  It cannot be $\mathrm{term}_A$, whose splits are all of
$x_2,\ldots,x_k$: that would give $g=0<f$.  So it is some
$\mathrm{dev}_j$ with $x_2=\cdots=x_j=1$, all drains before $j$ vanish,
and $\mathrm{dev}_j=1+d_j$.  As $\mathrm{dev}_2=1+d_2$ is also a height
and $f$ is the minimum, $d_j\le d_2$; the weights being nondecreasing,
$d_j=d_2$ and $f=1+d_2$ with $x_2=1$.
\end{proof}

\noindent
The balance solution \eqref{eq:balance} is the unique split vector making
\emph{all} $k$ heights equal, and it satisfies the first alternative.  The
second alternative is not vacuous, so the lemma cannot be strengthened to
$f=g$ on the closed box: at $k=4$ with $d=(\tfrac1{10},\tfrac15)$ the
splits $x=(1,\tfrac12,\tfrac15)$ make all three $A$-side heights equal to
$\tfrac{11}{10}$ while $\mathrm{term}_H=\tfrac{13}{10}$, and since
$f\le\mathrm{dev}_2\le1+d_2=\tfrac{11}{10}$ everywhere on the box, that
point maximizes $\min\{f,g\}$ with $f<g$.

Admissibility is what separates the two alternatives.  Over $2000$ random
nondecreasing weight vectors at depths $3$ to $7$, a linear program over
the splits returned the balanced value at every one of the $1336$
admissible vectors and the cap $1+d_2$ at every one of the $664$
inadmissible ones, without exception in either class.  We do not prove
this separation, but its consequence for the direction of error matters.
At an admissible weight vector the balanced value is, empirically,
exactly the split optimum, so \eqref{eq:tau} is not conservative there.
At an inadmissible one it \emph{overstates} what the box permits --- by
$\tfrac{123}{110}-\tfrac{11}{10}$ in the example above, where the balance
solution would need $x_2>1$.  This is why admissibility is not a
formality: every rung of Section~\ref{sec:certificates} is checked for it,
and the accompanying verifier refuses to certify a rung unless
$\tau\le1+d$ holds exactly.

\begin{theorem}[closed form]\label{thm:tau}
For every $k\ge3$,
\[
 \tau(k)\;=\;\max_{d\ \mathrm{admissible}}
 \frac{d\bigl(k-1-(k-2)\,d^{1/(k-2)}\bigr)+1}{1+d},
\]
and for every admissible $(k,d)$ and every $\eps>0$, suitable
$(p,\eta,N)$ in Lemma~\ref{lem:budget} give instances forcing discrepancy
$\ge\tau(k,d)-\eps$ with $D=1$; hence $\kappa\ge\tau(k,d)$.  Taking $d$
rational makes the whole certificate exactly checkable, which is how the
instances of Section~\ref{sec:certificates} are verified.
\end{theorem}

\begin{proof}
Combine \eqref{eq:tau} with Lemma~\ref{lem:amgm} and write $d=d_2$;
admissibility is precisely feasibility of the design.  For the
quantitative statement take $p=\eps/3$, $N=\lceil(\tau+1)/p\rceil$, and
$\eta=\eps/(3kN)$: then $\tau-p-(k(N-1)+1)\eta>\tau-\eps$ and
$Np-N(k-1)\eta>\tau$, so Lemma~\ref{lem:budget} applies.
\end{proof}

\section{The limit is a Lambert-$W$ constant}\label{sec:limit}

\begin{theorem}\label{thm:limit}
\[
 \lim_{k\to\infty}\tau(k)\;=\;1+\W(1/e),
\]
where $d^*=\W(1/e)$ is the unique root of $\ln d=-(1+d)$ in $(0,1)$.
Consequently $\kappa\;\ge\;1+\W(1/e)\;\approx\;1.2784645$.
\end{theorem}

\begin{proof}
Write $m=k-2$ and $\phi_m(d)=1+m\bigl(1-d^{1/m}\bigr)$, so that
$\tau(k,d)=[d\,\phi_m(d)+1]/(1+d)$.  For fixed $d\in(0,1)$ the quantity
$m(1-d^{1/m})=m(1-e^{\ln d/m})$ increases in $m$ to $-\ln d$ (apply
$e^x\ge1+x$ with $x=\ln d/m$), so
\[
 \tau(k,d)\ \uparrow\ \tau_\infty(d):=\frac{d(1-\ln d)+1}{1+d}
 \qquad(k\to\infty).
\]
Differentiating $\tau_\infty$, the numerator of the derivative is
\[
 (1+d)\frac{\partial}{\partial d}\bigl[d-d\ln d+1\bigr]
 -\bigl[d(1-\ln d)+1\bigr]
 \;=\;-\ln d-d-1,
\]
which decreases from $+\infty$ (as $d\downarrow0$) to $-2$ (at $d=1$); so
$\tau_\infty$ increases on $(0,d^*]$ and decreases on $[d^*,1]$, where
$\ln d^*=-(1+d^*)$, equivalently $d^*e^{d^*}=1/e$, i.e.\ $d^*=\W(1/e)$
\citep{CorlessGonnetHareJeffreyKnuth1996}.  Substituting
$\ln d^*=-(1+d^*)$,
\[
 \tau_\infty(d^*)=\frac{d^*(1+1+d^*)+1}{1+d^*}
 =\frac{(1+d^*)^2}{1+d^*}=1+d^*.
\]
\emph{Upper bound.}  For any $k$ and any admissible $d$ we have
$\tau(k,d)\le\tau_\infty(d)\le\tau_\infty(d^*)=1+d^*$, hence
$\tau(k)\le1+d^*$ for every $k$.

\emph{Lower bound.}  Fix $d\in(d^*,1)$.  Then
$\tau_\infty(d)\le1+d^*<1+d$, so $\tau(k,d)<1+d$ for \emph{every} $k$:
the pair $(k,d)$ is always admissible.  Since
$\tau(k,d)\uparrow\tau_\infty(d)$, we get
$\liminf_k\tau(k)\ge\tau_\infty(d)$, and letting $d\downarrow d^*$ (where
$\tau_\infty$ is continuous) gives $\liminf_k\tau(k)\ge1+d^*$.  The two
bounds give the limit.  Finally $\kappa\ge\tau(k,d)$ for every admissible
pair by Theorem~\ref{thm:tau}, so $\kappa\ge\tau_\infty(d)$ for every
$d>d^*$, and $\kappa\ge1+d^*$ follows.
\end{proof}

\begin{remark}
The identity $\tau_\infty=1+d^*$ says: at the optimum, the limiting chain
forces exactly one plus its own starting weight.  We do not know a direct
combinatorial explanation; it drops out of the stationarity condition.
\end{remark}

\section{Exact certificates}\label{sec:certificates}

All claims above are verified by machine in exact rational arithmetic.
Two independent toolchains are used: a certificate checker that
re-derives every inequality of Lemmas~\ref{lem:heights}
and~\ref{lem:budget} for a given $(k,d,p,\eta,N)$ from scratch using
exact fractions, and, for one concrete family member, an independent
mixed-integer feasibility check of the assembled instance.

\begin{center}
\begin{tabular}{@{}llllr@{}}
\toprule
depth & weights & $\tau$ (exact) & $\tau$ (dec.) & certified instance LB\\
\midrule
$k=3$ & $d=5/12$ & $239/204$ & $1.171569$ & $1.161469$\\
$k=4$ & $d=(9/25,3/5)$ & $103/85$ & $1.211765$ & $1.201665$\\
$k=5$ & $d=(1/3,10/21,7/10)$ & $904/735$ & $1.229932$ & $1.219832$\\
$k=7$ & near-geometric & $211933/170170$ & $1.245419$ & $1.235320$\\
$k=12$ & $r=22/25$ & (exact rational) & $1.261401$ & $1.251301$\\
$k=25$ & $r=473/500$ & (exact rational) & $1.270876$ & $1.269776$\\
$k=40$ & $r=967/1000$ & (exact rational) & $1.273843$ & $\mathbf{1.272743}$\\
\bottomrule
\end{tabular}
\end{center}

\noindent
Here $r$ denotes the common ratio of a rational geometric chain
$d_j=r^{k-j}$, which keeps every balance split \eqref{eq:balance} rational
($x_{j+1}=r$).  The certified instance lower bounds instantiate
Lemma~\ref{lem:budget} with $p=1/100$, $N=300$ for $k\le12$ and
$p=1/1000$, $N=3000$ for the deep rungs, $\eta=1/(10^4kN)$ throughout;
each is an explicit finite instance with $D=1$ and every column of
fractional support two.  The depth-$4$ member with $p=1/50$, $N=59$
($m=119$, $n=236$) was additionally checked by an exact mixed-integer
solver: infeasible at budget $7/6$, feasible at $1.19$.  This certifies
the promise made in the companion paper that deeper detectors beat
$4-2\sqrt2$: already $k=4$ reaches $103/85>4-2\sqrt2>7/6$.

\section{Beyond chains: what is true, and what fails}\label{sec:beyond}

Everything in this section is an exactly verified statement about finite
\emph{block games}, together with an honest account of why none of it
currently yields instances above the chain limit.  A block game fixes a
bank state (starting heights on designated rows), a budget $c$, and a
block of columns; all quantifiers over schedules are discharged by
exhaustive search in exact arithmetic.

\paragraph{Banked chains.}
If the helper row enters a chain block at height $h\ge0$, the analysis of
Sections~\ref{sec:gadget}--\ref{sec:optimize} goes through with
$\mathrm{term}_H$ increased by $h$, giving
\[
 \tau(k,h)=\max_{d\ \mathrm{adm.}}
 \frac{d\bigl(k-1+h-(k-2)d^{1/(k-2)}\bigr)+1}{1+d},
 \qquad
 \lim_k\tau(k,h)=1+\W(e^{h-1}),
\]
and on the compliant path the helper ends at exactly $\tau-1$.  Iterating
$h\mapsto\tau(\cdot,h)-1$ has fixed point $h^*=1/e$, suggesting a bound of
$1+1/e\approx1.3679$.  \emph{This is not a theorem}, and our experiments
show why the natural proof fails: sustaining $h$ across blocks requires a
guarantee over \emph{all} legal schedules, and a single legal gift of a
terminal column to the helper resets its bank from $h$ to $h-1$.  The
exists-a-schedule bookkeeping that suffices for one block does not
compose.

\paragraph{Super-threat blocks and shapers (exact).}
Two block-game facts, each verified by exhaustive exact search over all
schedules ($4^n$ plays, pruned):
(i) there is an explicit $14$-column block that, entered at bank state
$(1/2,9/20)$ on its two designated rows and with its detector column
denied, forces every schedule to a prefix height of exactly
$1.687669\ldots>3/2$, while admitting a legal escape play of value
$0.6931$ when the detector is granted;
(ii) there are \emph{cold shaper} blocks with no entry banks at all such
that, at every budget $c\le1.35$, every $c$-legal schedule exits with
some designated row at height $\ge0.3088$ --- verified at $c=1.35$, which
suffices for all smaller budgets because lowering $c$ only removes legal
schedules and so can only raise the minimum.  Further stage games, each
verified separately at its stated entry state, lift such banks in exact
steps: at budget $1.48$ the same block, entered at $0.31$ and again at
$0.43$, forces every legal schedule to exit at $0.43$ resp.\ $0.55$ (or
else to hand back a helper bank), the two stages having identical exit
margins $+0.1228$ because entry and threshold translate together.  A bank
of $0.65$ then fires: the denial of a detector column of share $49/50$
from that height is worth exactly $1.63$ on every play, above the budget
$1.60$ at which the corresponding stage guarantees were verified (and
firing is monotone in the bank, so any larger bank fires too).  We
emphasize that these are separate block games, each at its stated entry
state; whether they can be made to occur in sequence inside one instance
is precisely the composition question below.  At the level of components,
however, no budget we tested looks safe: for every budget up to $1.60$ we
found stage games whose guarantees, \emph{if composable}, would force a
violation beyond it, and the translation invariance just noted suggests
the pattern continues.

\paragraph{The composition barrier (measured).}
None of this composes.  In a fixed instance the scheduler plays all
blocks on shared rows, and three exactly measured mechanisms defeat every
assembly we tried.  \emph{Entry pollution}: stage guarantees proved at a
bank state are vacuous when the scheduler arrives elsewhere in the
$\pm c$ corridor, which it can, having parked rows anywhere within
budget.  \emph{Accumulator venting}: when the accumulator nears the
budget, one sacrificial receipt (a detector column granted to the
accumulator itself) vents it by nearly a full weight, and pre-parking the
pressured row makes the accompanying denial spike harmless.  In the
assembled $1273$-column machine below this is visible exactly: along the
certified schedule the accumulator's prefix height peaks at
$\tfrac{27999}{25000}$ and one column later stands at
$\tfrac{3999}{25000}$, a net relief of $\tfrac{24}{25}$, because that
detector column --- on which the accumulator carries weight $1$ --- is
granted to the accumulator itself.  \emph{No conjunctive pinning}:
hill-climbing over our block class, evaluating each candidate exhaustively
over all its legal schedules, produced no block forcing \emph{both}
designated rows to exit above useful thresholds; the best candidates leave
exact deficits of $-0.3242$, $-0.4892$ and $-0.5094$.  Exits are
disjunctive, so the scheduler always has a coordinate to wreck.  The
bottom line is
quantitative: the assembled machine built from components with
stage-level guarantees at $c=1.35$ admits a schedule with maximum prefix
discrepancy exactly $27999/25000=1.11996$ --- weaker than the plain
depth-$4$ chain instance, which forces $1.2017$.  (This is a witness, so
it bounds that machine's min-max from above; the accompanying code replays
the schedule in exact arithmetic.)  Block strength is unbounded;
composition is the constraint.

\paragraph{Conjecture and program.}
\begin{conjecture}\label{conj:finite}
$\kappa<\infty$.
\end{conjecture}
Two further data points support this.  First, a simple online rule
(relieve the taller support row unless doing so would sink it below a
fixed floor; otherwise relieve the other; otherwise dump on the row with
most headroom) already achieves maximum prefix discrepancy $1.11996$ on
our strongest assembled instance, and comes within $0.0040$ and $0.0021$ of
the exact optimum on the companion paper's $11\times15$ and $49\times72$
instances respectively; its worst case over adversarial search is below
$2.2$ on everything we could construct.  Second, the mechanisms
that defeat composition --- corridor slack, venting, disjunctive exits ---
are exactly the ingredients a potential-function upper-bound proof would
formalize: the scheduler maintains every row in a managed corridor and
can always afford one sacrificial relief before any accumulator bursts.
Adapting the earliest-deadline analysis of \citep{LiuReis2026} with
row-specific deadlines and such a potential is, in our view, the most
promising route to any finite upper bound; even $\kappa\le O(1)$ would be
the first such statement for machine-dependent weights.  In the other
direction, any instance forcing more than $1+\W(1/e)$ must defeat the
corridor mechanisms simultaneously on all rows, which our experiments
consistently failed to arrange.

\section{Bounded weight ratio}\label{sec:rho}

For applications the relevant parameter is often the dynamic range of the
weights rather than their absolute scale.  Define $\kappa(\rho)$ as
$\kappa$ restricted to instances with $d_{ij}\in[D/\rho,D]$.  Here
$\kappa(1)=1$: the upper bound is the common-weight theorem of Liu and
Reis, and Tijdeman's classical lower bound $1-1/(2m-2)-\delta$ supplies
the matching supremum \citep{Tijdeman1973,Meijer1973}.  Our constructions use tiny off-support
weights, but only through the bookkeeping of Lemma~\ref{lem:budget}:
setting $\eta=1/\rho$ there gives, for any admissible depth-$k$ point
with value $\tau$ and any $p$,
\[
 \kappa(\rho)\ \ge\
 \min\Bigl\{\tau-p-\tfrac{k(N-1)+1}{\rho},\ Np-\tfrac{N(k-1)}{\rho}\Bigr\},
\]
and choosing $N=\lceil2/p\rceil$, $p=2\sqrt{k/\rho}$ yields
$\kappa(\rho)\ge\tau(k,d)-4\sqrt{k/\rho}$ once $\rho\ge16k$.  Since
$\sup_{k,d}\tau(k,d)=1+\W(1/e)$, it follows that
$\kappa(\rho)\to1+\W(1/e)$ as $\rho\to\infty$; for a concrete data point,
$\rho\ge3.2\times10^4$ already gives $\kappa(\rho)>7/6$ via the depth-$4$
witness.  We make no attempt to optimize the rate, which is certainly far
from tight.  Determining the shape of $\kappa(\rho)$ between
$\kappa(1)=1$ and the machine-dependent regime --- mild heterogeneity
versus mixed fleets, in scheduling terms --- is open and seems the right
quantitative target for practice.

\section{Open problems}

\begin{enumerate}
\item Is $\kappa$ finite (Conjecture~\ref{conj:finite})?  Any finite
upper bound would be the first for machine-dependent prefix discrepancy.
\item Is $\kappa=1+\W(1/e)$?  The chain limit is the strongest
instance-level lower bound we know, all mechanisms beyond it having
failed to compose; but we see no conservation obstruction above it, and
the block-game results show local strength is not the bottleneck.
\item Is the balanced design optimal among all splits at every admissible
weight vector?  Lemma~\ref{lem:tight} gives only the dichotomy, and the
$2000$-vector separation reported after it is evidence, not proof.  Is the
chain in turn optimal among all single-block gadgets?  Our search
evidence (some $16{,}000$ mixed-integer evaluations over tree-shaped
punishment structures at depths $4$--$7$, spot-validated against
exhaustive exact search) found nothing better, but this is not a proof.
\item Determine $\kappa(\rho)$ asymptotics and, for practice, its profile
at moderate $\rho$.
\end{enumerate}

\section*{Reproducibility and disclosure}

The accompanying repository contains exact verifiers for every inequality
in Lemmas~\ref{lem:heights} and~\ref{lem:budget} at all seven table
rungs, including the deep geometric chains ($\texttt{Fraction}$-exact
throughout); the mixed-integer confirmation of the $119\times236$
instance; the exhaustive block-game verifiers behind
Section~\ref{sec:beyond} (branch-and-bound over all schedules in scaled
integer arithmetic); and the experiment logs, including documented retractions of intermediate
findings that deeper search or exact re-verification overturned.
Generative AI systems (GPT-5.6 Pro, Claude Opus 5, and Claude Fable 5)
assisted with exploration, computation, verification code, adversarial
review, and exposition.  They are not authors.  The author assumes
responsibility for independent verification, priority, attribution, and
the final text.

\bibliographystyle{abbrvnat}
\bibliography{references}

\end{document}
